\section{Invariant sections and anticanonical nonvanishing}
\label{sec:torus}

The ordinary theorem is now available on every no-forms base. We apply
it to a rational torus quotient in order to obtain a section with the
invariance required by Section~\ref{sec:structure}. The quotient
provides adjoint rank one for a geometric reason: a rational function
constant on a dense orbit is a function on the base. Smoothness of the
anticanonical metric provides positivity of that invariant line.
We first establish the required smooth variant of transfer, then
verify its hypotheses on the quotient.

\subsection{The smooth invariant transfer}
\label{sec:invariant-transfer}

On a torus quotient it is the space of invariant adjoint sections,
rather than the full section space, that is one-dimensional.
We need a different positivity input for that line, while the divisor
normalization and boundary estimates remain those already proved.

\begin{proposition}[Smooth invariant variant]\label{prop:invariant-transfer}
Use the varieties, maps, line bundles, and section of
Proposition~\ref{prop:transfer}, with \(D=0\).
Let a compact torus \(T\) act holomorphically and continuously on
\(X,\bar Y,Y\), with compatible holomorphic linearizations of the
line bundles. Require \(r,\bar\pi,\bar f\) to be equivariant, with
trivial action on \(S\). Require \(\rho\) and the \(J\)-valued adjoint
section \(\eta_Y=\pi^*s_0\) to be invariant, and require \(h_J,h_L\)
to be smooth, semipositive, and \(T\)-invariant.

Replace the ordinary rank hypothesis \eqref{eq:rank} by
\[
 \dim_{\C(S)}
 H^0\!\left(Y_\eta,
 \left(K_{Y/S}+\pi^*(J+jkL)\right)|_{Y_\eta}\right)^T=1
 \qquad(j=0,1,2,\ldots).
\]
Then all the conclusions of Proposition~\ref{prop:transfer} hold,
and the section on \(X\) produced by its lifting is \(T\)-invariant.
\end{proposition}

\begin{proof}
The normalized line \(A\) is given the trivial action.
In \eqref{eq:minima} only base divisors are used, so dividing the
invariant \(\rho\) by their rational section gives an invariant \(e\).
The sections \(u_j\) in \eqref{eq:moments} are invariant as well.

On a dense smooth-fibration open where the adjoint direct image is
locally free and its formation commutes with base change, Berndtsson's theorem
\cite[Theorem~1.2]{Berndtsson2009} gives its smooth \(L^2\) metric
semipositive curvature. Its hypotheses hold because the total space
\(Y\) is projective, hence K\"ahler, and each twisting metric
\(\pi^*(h_Jh_L^{jk})\) is smooth and semipositive.

Average the holomorphic group action on this vector bundle over
the compact torus. This gives a holomorphic projection onto its
invariant subbundle, which has rank one by assumption.
The fiber integral is invariant under the group, so the projection
is orthogonal. The metric on the invariant line is therefore its
quotient metric and is semipositive. Since \(u_j\) is its nonzero
local generator, \(-\log V_j\) has the required psh representative.

This is the only positivity step that changes. The fixed-open
estimate \eqref{eq:common-open}, the zeroth-moment extension
\eqref{eq:base-volume-weight}, the high-moment envelope
\eqref{eq:envelope}, and both bounds through codimension two
now apply exactly as in the ordinary proof.
They give the two asserted base metrics.
Every base section is invariant because the base action is trivial.
The rational lifting multiplies its pullback by \(e^a\), so remains
invariant. With \(D=0\), the support condition permits only
exceptional poles, and the final extension on \(X\) is an invariant
holomorphic section.
\end{proof}

The distinction between the two propositions is substantive.
The ordinary argument allows a singular volume metric because the
whole adjoint space has rank one and the relative Bergman theorem
applies. The invariant argument uses smooth coefficients to obtain
a positive quotient metric on an orthogonal holomorphic summand.

\subsection{Natural invariant nonvanishing}

\begin{theorem}[Natural invariant nonvanishing]\label{thm:invariant}
Let $Z$ be smooth connected projective over $\C$, with
$H^0(Z,\Omega_Z^q)=0$ for every $q>0$. Suppose that
$L=-K_Z$ has a smooth Hermitian metric of semipositive curvature.
Let a compact connected real torus $T$ act continuously on $Z$ by
holomorphic automorphisms, and give $L$ the natural linearization
induced by the differential of the action. Then
\[
                  H^0(Z,mL)^T\ne0
                  \quad\text{for some integer }m>0.
\]
\end{theorem}

\begin{proof}
If $Z$ is a point the natural action on its canonical line is trivial,
so the constant section proves the claim. Otherwise replace $T$ by
its image in $\operatorname{Aut}(Z)$. The no-forms hypothesis and
Hodge symmetry give $H^1(Z,\OO_Z)=0$. The algebraic-torus argument
in the proof of Theorem~\ref{thm:structure} applies to this action:
connectedness preserves the first Chern class of a very ample line,
the vanishing of $H^1(\OO_Z)$ makes that line invariant up to
isomorphism, and its projective representation embeds $T$ in a
compact diagonal torus. Its Zariski closure $G$ is an algebraic torus
acting on $Z$, and natural sections are $T$-invariant exactly when
they are $G$-invariant.

Average the given metric over $T$ by averaging local weights, using
the natural linearization. Differences of metric weights are global
functions, so these averages patch to a smooth invariant metric; its
curvature remains semipositive. Take $J=L$ and $D=0$. The canonical
section $1$ of $K_Z+L=\OO_Z$ is invariant for the natural action,
and its metric pairing is a smooth positive volume form of finite
total mass on compact $Z$.

If the general $G$-orbit has dimension zero, connectedness of $G$
makes its action trivial: the orbit map is constant on a dense open
and hence on $Z$. Theorem~\ref{thm:finite-volume}, with $D=0$ and
$J=L$, then gives the required section. Suppose henceforth that the
general orbit has dimension $r>0$.

Take the Rosenlicht rational quotient
\cite{Rosenlicht1956,BellGhiocaReichstein2017}. Its function field is
$\C(Z)^G$, and a dense invariant open has geometric fibers equal
to orbits. Because $G$ is connected, its general orbit is geometrically
integral; thus the invariant field is relatively algebraically closed
in $\C(Z)$ and the quotient has connected general fiber. Choose its
prepared equivariant model
\[
      Y\xrightarrow{r_Y}\bar Y\xrightarrow{\bar\pi}Z,
          \qquad f:Y\to S,\quad\bar f:\bar Y\to S,
\]
with trivial base action. The notation $r_Y$ for the resolution avoids
confusion with the orbit dimension $r$. The generic fiber has dimension
$r$ and a geometrically dense $G$-orbit. After shrinking a base open,
$f$ is smooth there, its fibers have a dense orbit, and $\pi$ agrees
birationally with the original quotient construction.

We construct the relative section required for invariant transfer.
Choose $r$ elements of the Lie algebra of $G$ whose fundamental vector
fields $v_1,\ldots,v_r$ are generically independent. Such a choice is
possible because the general orbit has dimension $r$. The vector
fields are invariant, since $G$ is commutative, and tangent to the
fibers of $f$. On the smooth-fibration open their wedge is a regular,
generically nonzero section of $-K_{Y/S}$. Tensor it with the pullback
of a rational frame of $-K_S$, regular after shrinking the base open.
This gives an invariant anticanonical section on $Y$ over that open.
The determinant of $d\pi$ is the natural morphism
\[
                     -K_Y\longrightarrow\pi^*(-K_Z).
\]
Its image is a nonzero invariant rational section of $\pi^*L$,
regular over the same open.

This rational section also defines a section $\rho$ of
$\bar\pi^*L$ on $\bar Y$ with no horizontal poles. Indeed,
$r_Y$ is an isomorphism at the generic point of each prime divisor
of normal $\bar Y$. Every horizontal prime meets the base open on
which the constructed section is regular. The birational identification
of function fields therefore gives nonnegative order there.

The generic-fiber adjoint spaces are finite-dimensional
$\C(S)$-vector spaces carrying algebraic representations of the
split torus $G_{\C(S)}$. They decompose into character spaces.
A character of $G$ is trivial on the Zariski-dense subgroup $T$
exactly when it is the zero character. Consequently their
$T$-fixed and $G$-fixed subspaces agree, not merely those of sections
on $Z$ itself.

For each $j\ge0$, on the generic fiber the section
\[
             (\pi^*1/f^*ds)\rho^j
\]
is a nonzero invariant section of
$K_{Y/S}+\pi^*((j+1)L)$. The form $\pi^*1$ includes the effective
birational Jacobian, and $\rho$ has no horizontal pole, so this is
regular. Any two nonzero invariant sections of this line have an
invariant rational function as their ratio. The equivariant birational
maps and the defining property of the quotient give
\[
                    \C(Y)^G=\C(Z)^G=\C(S),
\]
so this ratio belongs to $\C(S)$. Their
space has dimension one over $\C(S)$, for every $j\ge0$ including
$j=0$.

Proposition~\ref{prop:invariant-transfer} now applies with $k=1$, with powers
subsequently taken there to clear the vertical minima. It supplies a
line $A$ and effective divisor $D_S$ on $S$, with a bounded
semipositive metric on $A$ and a finite-volume semipositive metric on
$-K_S+D_S$. Holomorphic forms on $S$ pull back injectively to $Z$,
as in the last part of Section~\ref{sec:induction}; hence $S$ has no
positive-degree holomorphic forms. Theorem~\ref{thm:finite-volume}
gives a nonzero section of $aA+bD_S$ with $a>0,b\ge0$.
Invariant lifting produces a nonzero invariant section of a positive
power of $L$ on $Z$. Since $D=0$, every permitted pole upstairs is
exceptional over $Z$ and disappears on its normal pushdown.
When $S$ is a point the same conclusion is the point-base case of
Proposition~\ref{prop:invariant-transfer}. This completes the proof.
\end{proof}

\subsection{Descent to the original variety}

\begin{proof}[Proof of Theorem~\ref{thm:headline}]
Use Theorem~\ref{thm:structure} to choose the finite \'etale cover
$\nu:S\to X$ and the universal product
$\widetilde X=\C^a\times F\times Z$ with its lattice deck action.
The compact factor $Z$ satisfies the hypotheses of
Theorem~\ref{thm:invariant}; let $s_Z\ne0$ be an invariant section
of $K_Z^{-m}$ for some $m>0$. If $Z$ is a point, take $m=1$ and
$s_Z=1$.

Using the product identification \eqref{eq:cover-canonical}, form
\[
     \eta_0^{-m}\otimes\eta_F^{-m}\otimes s_Z
                \in H^0(\widetilde X,K_{\widetilde X}^{-m}),
\]
where pullbacks by the product projections are understood. The first
two frames are invariant under the deck lattice and the last factor
is invariant under the closure of its action. The product is therefore
invariant and descends to a nonzero holomorphic section of $K_S^{-m}$.
It is algebraic by projectivity. Finally Lemma~\ref{lem:etale-norm}
produces a nonzero section of $K_X^{-m\deg\nu}$. The exponent is
positive, as required.
\end{proof}
